\subsection{伯努利多项式}

考虑由下式定义的复合算子 U

\[
U (f (\mathrm{X})) = f (\mathrm{X} + 1) - f (\mathrm{X});
\]

可以写 \(U = e^{\mathrm{D}} - 1\) （VI, p.~270, 例 1）；这是一个 1 阶算子，且若令 \(U = DV\) ，则有 \(V^{-1} = \frac{\mathrm{D}}{e^{\mathrm{D}} - 1}\) 。\(n\) 次阿佩尔多项式（与算子 \(U\) 相应的）称为 \(n\) 次伯努利多项式，记作 \(\mathrm{B}_n(\mathrm{X})\) ；若令 \(b_n = \mathrm{B}_n(0)\) ，则有公式

\[
\mathrm{B} _ {n} (\mathrm{X}) = \sum_ {k = 0} ^ {n} \binom{n}{k} b _ {n - k} \mathrm{X} ^ {k}\tag{18}
\]

\[
\frac {\mathrm{Se} ^ {\mathrm{XS}}}{e ^ {\mathrm{S}} - 1} = \sum_ {n = 0} ^ {\infty} \frac {1}{n !} \mathrm{B} _ {n} (\mathrm{X}) \mathrm{S} ^ {n}\tag{19}
\]

特别地，有

\[
{\frac {\mathrm{S}}{e ^ {\mathrm{S}} - 1}} = \sum_ {n = 0} ^ {\infty} {\frac {1}{n !}} b _ {n} \mathrm{S} ^ {n}\tag{20}
\]

VI, p.~273 的公式 (7) 和 (9) 给出，对伯努利多项式，有关系

\[
\frac {d \mathrm{B} _ {n}}{d \mathrm{X}} = n \mathrm{B} _ {n - 1} (X)\tag{21}
\]

\[
\mathrm{B} _ {n} (\mathrm{X} + 1) - \mathrm{B} _ {n} (\mathrm{X}) = n \mathrm{X} ^ {n - 1}.\tag{22}
\]

特别地，有 \(\mathrm{B}_{n}(1)-\mathrm{B}_{n}(0)=0\) （对 n\textgreater1），由此并利用 (18)，就得到递推关系

\[
\sum_ {m = 0} ^ {n - 1} \binom {n} {m} b _ {m} = 0 \qquad (\text { for } n > 1)\tag{23}
\]

它使我们能逐步算出诸 \(b_{n}\) 。这些数显然都是有理数；由于可以写

\[
{\frac {S}{e ^ {S} - 1}} = - {\frac {S}{2}} + {\frac {S}{2}} {\frac {e ^ {S} + 1}{e ^ {S} - 1}}
\]

以及由于（VI, p.~272, 公式 (5)）

\[
\frac {e ^ {- S} + 1}{e ^ {- S} - 1} = - \frac {e ^ {S} + 1}{e ^ {S} - 1}
\]

可以看出，在 \(\frac{S}{2}\frac{e^{S}+1}{e^{S}-1}\) 的形式级数中，一切奇数次项的系数均为零；于是有

\[
b _ {0} = 1, \qquad b _ {1} = - \frac {1}{2}, \qquad b _ {2 n - 1} = 0 \quad \text { for } n > 1.\tag{24}
\]

有理数 \(b_{2n}\) （ \(n \geqslant 1\) ）称为伯努利数；我们将看到（VI, p.~288），\(b_{2n}\) 的符号是 \((-1)^{n-1}\) 的符号。公式 (23) 给出，对 n 的前几个值，

\[
\begin{array}{c} b _ {2} = \frac {1}{6}, \quad b _ {4} = - \frac {1}{3 0}, \quad b _ {6} = \frac {1}{4 2}, \quad b _ {8} = - \frac {1}{3 0}, \\ b _ {1 0} = \frac {5}{6 6}, \quad b _ {1 2} = - \frac {6 9 1}{2 7 3 0}, \quad b _ {1 4} = \frac {7}{6}, \\ b _ {1 6} = - \frac {3 6 1 7}{5 1 0}, \quad b _ {1 8} = \frac {4 3 8 6 7}{7 9 8}, \quad b _ {2 0} = - \frac {1 7 4 6 1 1}{3 3 0}, \quad b _ {2 2} = \frac {8 5 4 5 1 3}{1 3 8}, \\ b _ {2 4} = - \frac {2 3 6 3 6 4 0 9 1}{2 7 3 0}, \quad b _ {2 6} = \frac {8 5 5 3 1 0 3}{6}, \quad b _ {2 8} = - \frac {2 3 7 4 9 4 6 1 0 2 9}{8 7 0}. \end{array}
\]

注意 691、3617、43867 这几个分子是素数；其余的素因子分解为

\[
\begin{array}{c} 1 7 4 6 1 1 = 2 8 3 \times 6 1 7 \\ 8 5 4 5 1 3 = 1 1 \times 1 3 1 \times 5 9 3 \\ 2 3 6 3 6 4 0 9 1 = 1 0 3 \times 2 2 9 4 7 9 7 \\ 8 5 5 3 1 0 3 = 1 3 \times 6 5 7 9 3 1 \\ 2 3 7 4 9 4 6 1 0 2 9 = 7 \times 9 3 4 9 \times 3 6 2 9 0 3. \end{array}
\]

由此可导出前几个伯努利多项式的表达式

\[
\mathrm{B} _ {0} (\mathrm{X}) = 1, \quad \mathrm{B} _ {1} (\mathrm{X}) = \mathrm{X} - \frac {1}{2}, \quad \mathrm{B} _ {2} (\mathrm{X}) = \mathrm{X} ^ {2} - \mathrm{X} + \frac {1}{6},
\]

\[
\mathrm{B} _ {3} (\mathrm{X}) = \mathrm{X} ^ {3} - \frac {3}{2} \mathrm{X} ^ {2} + \frac {1}{2} \mathrm{X}, \quad \mathrm{B} _ {4} (\mathrm{X}) = \mathrm{X} ^ {4} - 2 \mathrm{X} ^ {3} + \mathrm{X} ^ {2} - \frac {1}{3 0}.
\]

